\documentclass[12pt]{article}
\usepackage[utf8]{inputenc}
\usepackage[T1]{fontenc}
\usepackage{hyperref}
\title{Kuthant Language Resources for K-12 Curriculum Development}
\author{Compiled by Claude Sonnet 5.5 (claude-sonnet-5-5) for Andr\'as Kornai}
\date{}
\begin{document}
\maketitle

\section{Introduction}
Kuthant (ISO 639-3 xut; Glottocode kuth1240), also called Gkuthaarn and by other spellings, is an extinct Paman language of the Cape York Peninsula, Queensland. Documentation is sparse. Some sources treat Garandi (Karundi) as a dialect of it, and the sources conflict on this. I found no revitalization programme or school materials.

\section{Primary Educational Resources}
None found.

\section{Linguistic and Reference Materials}
\begin{itemize}
\item Wikipedia, ``Gkuthaarn language'': overview and variant names. \url{https://en.wikipedia.org/wiki/Gkuthaarn_language}
\item Wikipedia, ``Garandi language'': discusses the relationship to Gkuthaarn. \url{https://en.wikipedia.org/wiki/Garandi_language}
\item Glottolog reference entries for Sidney H. Ray (1907), a comparative vocabulary of Cape York languages. I did not confirm whether it contains Kuthant data. \url{https://glottolog.org/resource/reference/id/472081}
\item AIATSIS manuscript finding aid MS 932, listing Kuthant grammatical notes among other material. \url{https://aiatsis.gov.au/sites/default/files/catalogue_resources/MS932.PDF}
\item Spanish Wikipedia mirror, ``Idioma kuthant'': a few words attributed to a secondhand 19th-century source; unverified. \url{https://wiki.undernet.uy/content/wikipedia_es_all_maxi_2025-08/Idioma_kuthant}
\end{itemize}
A selected bibliography of the language compiled by Lisa Marcussen is held in the AIATSIS Library according to search summaries; I did not find a URL.

\section{Religious and Community Resources}
None found in search results.

\section{Digital and Audio Resources}
No audio or digital resources were found. AIATSIS (AUSTLANG) lists the language with ``potential data'' status.

\section{Curriculum Development Context}
Very little is documented. Any programme should be led by descendant communities in the Gulf of Carpentaria area, and would begin by consulting AIATSIS holdings with the community's permission. The documentation is probably too thin for teaching the language itself. I did not identify a community organization.

\end{document}